\title{Direct Preference Optimization: Your Language Model is Secretly a Reward Model}

\begin{abstract}
Although large-scale unsupervised language models (LMs) acquire extensive world knowledge and exhibit some reasoning capabilities, directing their outputs with precision remains challenging due to the nature of unsupervised training. Current strategies to enhance model controllability typically involve gathering human judgments comparing the quality of different model responses and subsequently adapting the LM to match these preferences, often through reinforcement learning from human feedback (RLHF). However, the RLHF process is both intricate and prone to instability; it involves first training a reward model to capture human preferences, and then using reinforcement learning to update the LM so as to maximize the predicted reward while avoiding excessive deviation from the original model's behavior. In this work, we propose a novel reward model parameterization within RLHF that permits closed-form derivation of the optimal policy, thereby allowing the RLHF objective to be addressed via a straightforward classification loss. This approach, named \textit{Direct Preference Optimization} (DPO), is robust, effective, and computationally efficient. It removes the necessity for LM sampling during fine-tuning and diminishes the dependence on extensive hyperparameter searches. Empirically, we demonstrate that DPO fine-tunes LMs to better match human preferences as effectively as, or even surpassing, alternative methods. Specifically, DPO-based fine-tuning offers superior sentiment control over PPO-based RLHF, and delivers equal or higher response quality for summarization and single-turn dialogue tasks, while greatly simplifying implementation and training requirements.
\end{abstract}

\section{Introduction}
Large-scale unsupervised language models (LMs) trained on expansive datasets have demonstrated remarkable emergent abilities~\citep{chowdhery2022palm, brown2020language, touvron2023llama, bubeck2023sparks}. Despite these advances, such models are built from human-generated data that reflects a vast spectrum of intentions, priorities, and expertise levels. Many of these underlying motives or competencies may be unsuitable for reproduction in AI systems. For example, it is beneficial for an AI coding assistant to \textit{recognize} prevalent programming errors so it can correct them. Yet, during code generation, our aim is for the model to emphasize the (potentially uncommon) high-level coding proficiency found within its training corpus. Likewise, it may be useful for a language model to \textit{know} about a widespread misconception shared by half of the population, but it would be undesirable for the model to reproduce this misconception in half its answers when asked. Hence, isolating and enhancing the model’s \emph{preferred outputs and conduct}—as opposed to passively reflecting its vast repository of \textit{knowledge and competencies}—is fundamental to creating AI systems that are secure, effective, and manageable~\citep{ouyang2022training}. Although prevailing approaches use reinforcement learning (RL) to align LMs with human preferences, our findings indicate that the traditional RL-derived objective can in fact be directly optimized using a straightforward binary cross-entropy loss, greatly streamlining the preference learning workflow.

Generally, existing approaches impart target behaviors to language models by utilizing handpicked collections of human preference data, which exemplify actions deemed safe and useful by people. This process of learning from preferences is conducted following the initial phase of unsupervised pre-training on vast textual resources. While the most direct method to convey preferences involves supervised fine-tuning on exemplary human-crafted responses, the approaches with the greatest impact thus far rely on reinforcement learning from human or AI-generated feedback (RLHF/RLAIF; \citep{christiano2017deep, bai2022constitutional}). RLHF methodology involves first building a reward model trained on preference-annotated data, after which RL is applied to shape the LM policy to generate outputs that maximize this reward, all while constraining divergence from the original pre-trained policy. Although RLHF achieves LMs with notable interactive and coding skills, it adds substantial complexity to the model training process compared to standard supervised methods, requiring the training of several language models and repeated sampling from the model during optimization, thereby imposing considerable computational burdens.

In this work, we demonstrate a method for optimizing a language model to align with human preferences without relying on explicit reward modeling or reinforcement learning techniques. We introduce \textit{{Direct Preference Optimization} (DPO)}, an approach that implicitly targets the same goal as traditional RLHF methods—maximizing reward subject to a KL-divergence constraint—yet remains straightforward to implement and easy to train. Conceptually, DPO updates serve to boost the relative log likelihood of preferred over non-preferred responses, while incorporating a dynamic, example-specific importance weighting scheme that mitigates the model collapse seen with simpler probability ratio approaches. As with prior algorithms, DPO leverages a theoretical preference framework (such as the Bradley-Terry model; \cite{bradley1952rankanalysis}) for quantifying the alignment of a reward function to observed preference data. The key distinction is that, rather than using this preference model to define a loss for training a separate reward model and then updating the policy to optimize that reward model, DPO uses a variable transformation to formulate the preference loss directly as a function of the policy itself. This enables direct optimization of the policy based on human preference data using a straightforward binary cross-entropy loss, yielding the optimal policy relative to an implicit reward shaped by the preferences.

Our principal contribution is Direct Preference Optimization (DPO), a reward-model-free and RL-free technique for preference-driven training of language models. Empirical results show that DPO matches or outperforms existing methods, including RLHF approaches based on PPO, across tasks such as sentiment control, summarization, and conversation, using models with up to 6B parameters.

\section{Related Work}

As self-supervised language models have grown in size, they have demonstrated the capacity to handle certain tasks in a zero-shot setting \citep{radford2019language} or through the use of few-shot prompts \citep{gpt3,megatron,chowdhery2022palm}. Nonetheless, by fine-tuning these models on curated datasets containing instructions and corresponding human-generated responses, their effectiveness on downstream applications and alignment with users’ expectations can be greatly enhanced \citep{mishra-etal-2022-cross,sanh2022multitask,chung2022scaling,thoppilan2022lamda}. This process, often termed ‘instruction-tuning,’ has proven to facilitate the ability of LLMs to extrapolate to previously unseen instructions, thereby boosting their overall practicality \citep{chung2022scaling}. Although instruction-tuning has yielded considerable progress, it is generally less resource-intensive to gather \textit{relative} human assessments of response quality compared to assembling expert-level exemplars. Consequently, recent advancements have employed human preference data to further fine-tune LLMs, leading to better results across tasks such as translation \citep{kreutzer-etal-2018-reliability}, summarization \citep{stiennon2022learning,ziegler2020finetuning}, story-telling \citep{ziegler2020finetuning}, and instruction-following \citep{ouyang2022training,ramamurthy2023is}. The typical procedure involves first adapting a neural reward model using a preference model like the Bradley-Terry formulation \citep{bradley1952rankanalysis} to fit the collected pairwise preferences, and subsequently refining the language model via reinforcement learning strategies—commonly employing REINFORCE \citep{williams1992reinforce}, proximal policy optimization (PPO; \cite{schulman2017proximal}), or related techniques \citep{ramamurthy2023is}—to maximize this reward. In a related approach, LLMs already adjusted for instruction following via human input are further exploited to generate synthetic preference data that target specific attributes (e.g., safety or harmlessness) with only minimal human intervention, guided by text-based rubrics for the model’s output annotations \citep{bai2022constitutional}. These methodologies reflect a synthesis of two previously distinct areas: the application of reinforcement learning to language model training for diverse tasks~\citep{Ranzato2015SequenceLT,paulus2018a,wu2018learning}, and general strategies for leveraging human preference data \citep{christiano2017deep,kupcsik2018learning}. Despite the apparent benefits of utilizing comparative human judgments, the process of refining large language models through reinforcement learning continues to pose considerable technical difficulties. This work introduces a principled framework for optimizing over relative preferences that circumvents the need for reinforcement learning.

Beyond the realm of language modeling, research has explored the problem of learning policies from preferences in both the bandit and reinforcement learning domains, giving rise to a variety of methodological innovations. One notable example is contextual dueling bandits (CDB; \cite{yue2012karmed,dudik2015contextual}), where policies are learned based on pairwise preferences or action rankings rather than scalar rewards. In this context, because direct reward feedback is unavailable, theoretical work replaces the concept of an optimal policy with that of a \textit{von Neumann winner}—defined as a policy whose expected likelihood of outperforming any alternative policy is at least 50\% \citep{dudik2015contextual}. Despite these similarities, a key distinction exists: preference information in CDBs is obtained in an online fashion, whereas in human preference learning, the training process relies on a static, offline dataset composed of preference-annotated action pairs \citep{yan2022human}. An analogous research thread, \textit{preference-based RL} (PbRL), investigates learning from binary preferences that are sampled via an unobserved 'scoring' mechanism, not explicit rewards \citep{BusaFekete2014,ruiz2023dueling}. PbRL methods encompass a range of algorithms, some of which support leveraging off-policy preference data, but these approaches typically involve first learning an explicit model of the underlying reward function before conducting policy optimization \citep{jain2013learning,BusaFekete2014,christiano2017deep,sadigh2017active,kupcsik2018learning}. By contrast, we propose a one-stage approach that optimizes policies to adhere to preference information directly, without separate reward estimation.

\section{Preliminaries}\label{section:prelims}

We summarize the RLHF pipeline as described by \citeauthor{ziegler2020finetuning}, further extended in works such as \citep{stiennon2022learning, bai2022training, ouyang2022training}. The standard pipeline comprises three sequential steps: 1) supervised fine-tuning (SFT); 2) acquisition of preference data and reward model training; and 3) policy optimization through reinforcement learning.

\textbf{SFT}: The RLHF process usually initiates by adapting a large language model, pre-trained on general corpora, to specific tasks of interest—such as dialogue or summarization—by conducting supervised learning on curated datasets, thereby obtaining $\pi^\text{SFT}$.

\textbf{Reward Modelling Phase}: In the next step, prompts $x$ are fed into the SFT model to generate response pairs $(y_1, y_2)\sim \pi^\text{SFT}(y \mid x)$. Human annotators are then tasked with expressing a preference for one response over the other, represented as $y_w\succ y_l \mid x$, where $y_w$ is the chosen response and $y_l$ is the less preferred alternative from the pair $(y_1, y_2)$. These preferences are presumed to be induced by a latent reward function $r^*(y, x)$, to which we have no direct access. To model these human choices, various approaches are employed, with the Bradley-Terry (BT) model \cite{bradley1952rankanalysis} being especially prevalent (though more general models such as Plackett-Luce \citep{plackett1975analysis, luce2012individual} can be utilized when multiple ranked completions are available). The BT model formulates the observed preference distribution $p^*$ as follows:

\begin{equation}\label{eq:bradley-terry}
    p^*(y_1\succ y_2 \mid x)=\frac{\exp\left(r^*(x, y_1)\right)}{\exp\left(r^*(x, y_1)\right) + \exp\left(r^*(x, y_2)\right)}.
\end{equation}

Given a fixed dataset of pairwise preferences, $\mathcal{D}=\bigl\{x^{(i)}, y_w^{(i)}, y_l^{(i)}\bigr\}_{i=1}^N$, drawn from $p^*$, it is common to introduce a reward model $r_{\phi}(x, y)$ and train its parameters via maximum likelihood estimation. When treating the task as binary classification, the corresponding negative log-likelihood loss can be written as:

\begin{equation}\label{eq:reward_model}
    \mathcal{L}_R(r_{\phi}, \mathcal{D}) = -\mathbb{E}_{(x, y_w, y_l)\sim \mathcal{D}}\bigl[\log \sigma(r_{\phi}(x, y_w)- r_{\phi}(x, y_l))\bigr]
\end{equation}

Here, $\sigma$ denotes the logistic sigmoid function. For language modeling applications, the reward model $r_{\phi}(x, y)$ is usually built upon the SFT model $\pi^\text{SFT}(y \mid x)$, appending a linear head to the topmost transformer layer to yield a scalar output as the reward signal \cite{ziegler2020finetuning}. To stabilize the reward estimation and reduce its variance, prior research often normalizes reward outputs to enforce $\mathbb{E}_{x,y\sim \mathcal{D}}\left[r_\phi(x, y)\right] = 0$ for each $x$.

\textbf{RL Fine-Tuning Phase}: In the RL fine-tuning step, this reward model is used to guide learning of the language model. Consistent with previous approaches~\citep{jaques2017sequence, jaques2020human}, the optimization goal is defined as

\begin{equation}\label{eq:RL}
\max_{\pi_{\theta}}  \mathbb{E}_{x\sim \mathcal{D}, y\sim \pi_{\theta}(y \mid x)}\bigl[r_{\phi}(x, y)\bigr] - \beta\mathbb{D}_{\textrm{KL}}\bigl[\pi_{\theta}(y\mid x)\mid \mid \pi_\text{ref}(y\mid x)\bigr],
\end{equation}

where the hyperparameter $\beta$ regulates the divergence from the reference policy $\pi_\text{ref}$, typically chosen as the initial SFT model $\pi^\text{SFT}$. The policy $\pi_\theta$ is generally initialized from $\pi^\text{SFT}$ as well. This KL regularization is crucial to constrain the model's output space to regions where the reward model remains accurate and to maintain response diversity, avoiding mode collapse to a small set of high-reward completions. Since generating sequences in natural language is inherently discrete, the objective cannot be directly differentiated and instead requires reinforcement learning techniques for optimization. The prevailing method \citep{ziegler2020finetuning, stiennon2022learning, bai2022training, ouyang2022training} involves defining the adjusted reward as ${r(x, y) = r_{\phi}(x, y) -\beta (\log \pi_{\theta}(y\mid x) - \log \pi_\text{ref}(y\mid x))}$ and optimizing it with the PPO algorithm \cite{schulman2017proximal}.

\section{Direct Preference Optimization}\label{sec:DPO}

Driven by the complexities encountered with reinforcement learning for fine-tuning large language models, we set out to develop a more streamlined method for optimizing policies based directly on preferences. Unlike traditional RLHF strategies that separate reward modeling and subsequent RL-based policy training, our method introduces a specific reward parameterization, allowing derivation of the corresponding optimal policy in closed form, thus removing the necessity of an explicit RL process. As detailed below, our central idea is to exploit an analytical mapping from reward functions to their corresponding optimal policies, making it possible to reformulate the loss—originally over reward models—into a loss function defined on policies instead. This variable substitution approach obviates the need to train an isolated reward network, yet remains consistent with established preference modeling frameworks, such as the Bradley-Terry formulation. Effectively, this results in a

\textbf{DPO Objective Derivation.} We begin with the RL objective given in Eq.~\ref{eq:RL} as in earlier works, considering a general reward function $r$. As established by previous literature~\citep{peters2007reinforcement, peng2019advantage, korbak2022reinforcement, go2023aligning}, the solution that maximizes the KL-regularized reward in Eq.~\ref{eq:RL} assumes the following structure:

\begin{equation}\label{eq:op_policy}
    \pi_r(y\mid x) = \frac{1}{Z(x)}\pi_\text{ref}(y\mid x)\exp\left(\frac{1}{\beta}r(x, y)\right),
\end{equation}

where the normalization constant $Z(x) =\sum_{y}\pi_\text{ref}(y\mid x)\exp\left(\frac{1}{\beta}r(x, y)\right)$ ensures the distribution sums to one. The full derivation is provided in Appendix \ref{app:derivation1}. Even if we employ the MLE-estimated reward function $r_\phi$ in place of the true reward $r^*$, evaluating the partition function $Z(x)$ remains computationally prohibitive~\citep{korbak2022reinforcement, go2023aligning}, thereby limiting the tractability of this parameterization. However, Eq.~\ref{eq:op_policy} can be manipulated to express $r(x, y)$ in terms of $\pi_r$, $\pi_\text{ref}$, and $Z(x)$. Taking the logarithm of both sides of Eq.~\ref{eq:op_policy} and rearranging yields:

\begin{equation}\label{eq:main_eq}
    r(x,y) =\beta \log \frac{\pi_r(y\mid x)}{\pi_\text{ref}(y\mid x)} + \beta \log Z(x).
\end{equation}

This reformulation can be applied to the true reward function $r^*$ and its associated optimal policy $\pi^*$. Notably, in the Bradley-Terry model, only the difference in rewards for two responses is relevant, that is, ${p^*(y_1 \succ y_2 \mid x) = \sigma(r^*(x, y_1) - r^*(x, y_2))}$. Replacing $r^*(x,y)$ using Eq.~\ref{eq:main_eq} within the preference model Eq.~\ref{eq:bradley-terry} eliminates the partition function, allowing the human preference probability to be written solely in terms of the optimal policy $\pi^*$ and the reference policy $\pi_\text{ref}$. Therefore, under the Bradley-Terry preference structure, the optimal RLHF policy $\pi^*$ fulfills the following relationship:

\begin{equation}\label{eq:objective}
    p^*(y_1\succ y_2 \mid x)=\frac{1}{1 + \exp\left(\beta \log \frac{\pi^*(y_2\mid x)}{\pi_\text{ref}(y_2\mid x)} - \beta \log \frac{\pi^*(y_1\mid x)}{\pi_\text{ref}(y_1\mid x)}\right)}
\end{equation}

A detailed derivation appears in Appendix~\ref{app:derivation2}. Although Eq.~\ref{eq:objective} utilizes the Bradley-Terry model, parallel derivations can be constructed for more general cases, such as the Plackett-Luce models~\citep{plackett1975analysis, luce2012individual}, as elaborated in Appendix~\ref{app:plackett_luce_models}.

Given the expression for human preference likelihood as a function of the optimal policy, we can set up a maximum likelihood estimation objective over a parameterized policy $\pi_\theta$. Mirroring the approach used in reward model learning (see Eq.~\ref{eq:reward_model}), the objective for learning the policy can be stated as:

\begin{equation}\label{eq:optimum_model}
    \mathcal{L}_\text{DPO}(\pi_{\theta}; \pi_\text{ref}) = -\mathbb{E}_{(x, y_w, y_l)\sim \mathcal{D}}\left[\log \sigma \left(\beta \log \frac{\pi_{\theta}(y_w\mid x)}{\pi_\text{ref}(y_w\mid x)} - \beta \log \frac{\pi_{\theta}(y_l\mid

In this formulation, we infer an implicit reward via a different parameterization, whereby the optimal policy is directly represented as $\pi_\theta$. Furthermore, because our approach aligns with fitting a reparameterized Bradley-Terry model, it inherits several theoretical advantages, notably consistency under appropriate assumptions about the distribution of preference data \cite{bong2022generalized}. Additional discussion of these theoretical aspects and the relationship of DPO to prior work is presented in Section~\ref{sec:theory}.

\textbf{Mechanistic interpretation of the DPO update.} To gain insight into how DPO functions, it is informative to examine the gradient of the objective $\mathcal{L}_\text{DPO}$. The gradient with respect to the parameters $\theta$ is given by:

\begin{multline*}\label{eq:gradient}
    \nabla_\theta \mathcal{L}_\text{DPO}(\pi_\theta;\pi_\text{ref}) = \\ -\beta\mathbb{E}_{(x, y_w, y_l) \sim \mathcal{D}} \bigg[\underbrace{\sigma(\hat{r}_\theta(x, y_l) - \hat{r}_\theta (x, y_w))}_\text{emphasis on higher error in reward predictions}\bigg[\underbrace{\nabla_\theta\log \pi(y_w \mid x)}_\text{promote preferred output} - \underbrace{\nabla_\theta\log\pi(y_l \mid x)}_\text{discourage dispreferred output}\bigg]\bigg],
\end{multline*}

with the implicit reward computed as $\hat{r}_\theta(x, y) = \beta \log \frac{\pi_\theta(y \mid x)}{\pi_\text{ref}(y \mid x)}$, defined in terms of both the current policy $\pi_\theta$ and the reference $\pi_\text{ref}$ (see Section~\ref{sec:theory} for details). Intuitively, the gradient encourages increasing the probability assigned to the more preferred completions $y_w$, while reducing the probability of less favored completions $y_l$. Crucially, each training instance is weighted by the degree to which the current implicit reward model $\hat{r}_\theta$ overestimates the value of the less preferred completion, with this weighting modulated by $\beta$. This effectively measures the severity of reward mis-ranking and is influenced by the strength of the KL constraint. Our empirical findings underscore the significance of this weighting: omitting it (i.e., using a naïve variant that does not include the weighting coefficient) leads to significant degeneration in the language model’s output (refer to Appendix Table~\ref{tab:unlikelihood_generations}).

\textbf{DPO overview.} The standard DPO workflow involves the following steps: 1) For each prompt $x$, generate completions $y_1, y_2 \sim \pi_\text{ref}(\cdot \mid x)$, then annotate these with human preferences to assemble an offline preference dataset $\mathcal{D} = \{x^{(i)}, y_w^{(i)}, y_l^{(i)}\}_{i=1}^N$; 2) train the language model $\pi_\theta$ by minimizing the DPO loss $\mathcal{L}_\text{DPO}$, conditioned on the fixed $\pi_\text{ref}$, dataset $\mathcal{D}$, and chosen $\beta$. 
In practical settings, it is often desirable to make use of available public preference datasets, rather than collecting new samples and human preference labels. Typically, such datasets are derived using $\pi^\text{SFT}$; therefore, if $\pi^\text{SFT}$ exists, we set $\pi_\text{ref} = \pi^\text{SFT}$. If, however, $\pi^\text{SFT}$ is unavailable, $\pi_\text{ref}$ is instead initialized via maximum likelihood estimation over preferred completions $(x, y_w)$, specifically, $\pi_\text{ref} = \argmax_{\pi}\mathbb{E}_{x, y_w \sim \mathcal{D}}\left[\log \pi(y_w \mid x)\right]$. This initialization serves to bridge the gap between the ideal reference distribution (which is inaccessible) and the practical $\pi_\text{ref}$ utilized within DPO. Details concerning implementation and chosen hyperparameters can be found in Appendix~\ref{app:implementation}.

\section{Theoretical Analysis of DPO}
In this section, we deepen the analysis of the DPO approach, offering both a theoretical justification and insights that explain DPO’s benefits compared to actor-critic based RLHF methods such as PPO~\cite{schulman2017proximal}.

\label{sec:theory}

\subsection{Your Language Model Is Secretly a Reward Model} 
DPO accomplishes preference optimization without learning an explicit reward model or carrying out reinforcement learning by using a single maximum likelihood objective. Observe that the DPO loss (see Eq. \ref{eq:main_eq}) corresponds to the Bradley-Terry model with rewards defined as $r^*(x, y) = \beta \log\frac{\pi^*_\theta(y \mid x)}{\pi_\text{ref}(y \mid x)}$, where the model $\pi_\theta$ is trained analogously to the reward model in Eq. \ref{eq:reward_model} but under a suitable change of variables. In the following, we construct the theory underlying this reparameterization, demonstrating that it places no additional restrictions on the reward model class and that it is sufficient for recovering the optimal policy exactly. We start by formalizing an equivalence relation on reward functions.

\begin{definition}
Two reward functions $r(x, y)$ and $r'(x, y)$ are said to be equivalent if and only if ${r(x, y)-r'(x, y) = f(x)}$ for some function $f$.
\end{definition}

This equivalence clearly satisfies the requirements of an equivalence relation, thereby dividing the space of reward functions into distinct equivalence classes. Based on this definition, we can establish the following lemmas:

\begin{lemma}\label{lemma:same_prefrence}
Within the framework of Plackett-Luce, and by extension Bradley-Terry, preference models, any two reward functions that belong to the same equivalence class will give rise to identical preference distributions.
\end{lemma}

\begin{lemma}\label{lemma:same_policy}
Reward functions that are members of the same equivalence class result in the same optimal policy in the context of the constrained RL problem.
\end{lemma}

As the arguments are elementary, we postpone detailed proofs until Appendix \ref{app:lemma1}. The first lemma points to the well-recognized problem of under-specification associated with the Plackett-Luce model family \cite{plackett1975analysis}. Owing to this ambiguity, it is common practice to enforce supplementary identifiability conditions in order to provide meaningful guarantees for MLE estimators described by Eq.~\ref{eq:reward_model} \cite{bong2022generalized}. The second lemma asserts that within a given equivalence class, all reward functions induce an identical optimal policy. Thus, for our purposes, it suffices to identify any representative reward function from the target optimal equivalence class. We establish the following theorem, with its proof provided in Appendix~\ref{app:thm1}:

\begin{theorem}\label{thm:main}
    Assuming certain mild conditions, each reward class compatible with the Plackett-Luce family (including Bradley-Terry) admits a representation of the form ${r(x, y) = \beta \log \frac{\pi(y\mid x)}{\pi_\text{ref}(y\mid x)}}$, where $\pi(y\mid x)$ is some learned model and $\pi_\text{ref}(y \mid x)$ is a fixed reference model.
\end{theorem}

\begin{sproof}
    Take an arbitrary reward function $r(x, y)$ which determines a unique optimal model $\pi_r(y \mid x)$ given by Eq.~\ref{eq:op_policy}. Our aim is to demonstrate that a reward function from the equivalence class of $r$ admits the proposed reparameterization. Define the projection operator $f$ by 
\begin{equation}
    f(r; \pi_\text{ref}, \beta)(x, y) = r(x, y) - \beta\log\sum_{y}\pi_\text{ref}(y\mid x)\exp\left(\frac{1}{\beta}r(x, y)\right)
\end{equation}
Here, $f$ applies a normalization by subtracting the logarithm of the partition function with respect to the reference model $\pi_\text{ref}$, so the normalization term depends only on the context $x$. Consequently, $f(r; \pi_\text{ref}, \beta)(x, y)$ also lies within the equivalence class of $r(x, y)$. Substituting $r$ using the right-hand side of Eq.~\ref{eq:main_eq} (which is valid for any reward), it follows that $f(r; \pi_\text{ref}, \beta)(x, y) = \beta \log \frac{\pi_r(y\mid x)}{\pi_\text{ref}(y\mid x)}$. Therefore, this projection recovers a function in the equivalence class of $r$ in the required parametric form, ensuring no generality is lost by adopting the stated reward model.
\end{sproof}

Theorem~\ref{thm:main} can also be interpreted as uniquely identifying the specific reward function within each equivalence class that is chosen by the DPO reparameterization—namely, the reward function that fulfills:

\begin{equation}\label{eq:lag_p}
     \sum_{y}\underbrace{\pi_\text{ref}(y\mid x)\exp\left(\frac{1}{\beta}r(x, y)\right)}_{=\pi(y\mid x)\text{, using Thm.~\ref{thm:main} reparam.}} = 1,
\end{equation}

which ensures that $\pi(y\mid x)$ forms a legitimate probability distribution, where all probabilities are non-negative and sum to unity. Additionally, based on Eq.~\ref{eq:op_policy}, it becomes clear that Eq.~\ref{eq:lag_p} represents the partition function corresponding to the optimal policy determined by the reward function $r(x, y)$. The primary contribution of the DPO method lies in its ability to introduce constraints on the otherwise underdetermined Plackett-Luce family of preference models (including Bradley-Terry), such that the range of expressible reward models is preserved while simultaneously ensuring that the optimal policy described in Eq.~\ref{eq:op_policy} remains computationally manageable for all input prompts $x$.

\subsection{Instability of Actor-Critic Algorithms} 

Our framework can further elucidate sources of instability in conventional actor-critic approaches commonly applied in RLHF, such as PPO. In keeping with the RLHF process, we examine the RL fine-tuning phase described in Section \ref{section:prelims}. By mapping this scenario onto the control as inference paradigm \cite{levine2018reinforcement} for the constrained reinforcement learning formulation in \ref{eq:RL}, we postulate a parameterized policy $\pi_{\theta}(y\mid x)$ and seek to minimize $\mathbb{D}_{\text{KL}}[\pi_{\theta}(y|x) \mid \mid \pi^*(y\mid x)]$, where $\pi^*$ denotes the optimal policy generated by the reward $r_{\phi}(y, x)$ as defined in Eq.~\ref{eq:optimum_model}. Through straightforward manipulations, this setup leads to the following optimization problem:

\begin{equation}\label{eq:AC}
    \max_{\pi_{\theta}}\mathbb{E}_{\pi_{\theta}(y\mid x)}\bigg[\underbrace{r_{\phi}(x, y) -\beta\log\sum_{y}\pi_\text{ref}(y\mid x)\exp\left(\frac{1}{\beta}r_{\phi}(x, y)\right)}_{f(r_{\phi}, \pi_\text{ref}, \beta)} - \underbrace{\beta\log\frac{\pi_{\theta}(y\mid x)}{\pi_\text{ref}(y\mid x)}}_{\text{KL}}\bigg]
\end{equation}

This objective aligns with that used in previous research \citep{ziegler2020finetuning, stiennon2022learning, bai2022training, ouyang2022training}, where the DPO-equivalent reward corresponding to the reward class $r_{\phi}$ is optimized. In this framework, the normalization component within $f(r_{\phi}, \pi_\text{ref}, \beta)$ can be viewed as representing the soft value function of the reference policy $\pi_\text{ref}$. Although this term does not influence the location of the optimal solution, omitting it can significantly increase the variance of the policy gradient estimator, thereby destabilizing training. It is possible to account for this normalization through a learned value function, though this approach may itself present optimization challenges. A different strategy employed by earlier studies involves normalizing the reward using a human completion baseline, which serves as a one-sample Monte Carlo estimate of the normalization constant. The DPO parameterization, however, produces a reward function that eliminates the need for any such baseline.

\section{Experiments}
We empirically investigate the efficacy of DPO for preference-based policy training in this section. We first consider a controlled text generation task to examine how DPO balances reward maximization against KL-divergence minimization with respect to the reference policy, benchmarking its sample efficiency against established preference learning approaches such as PPO. Subsequently, we assess DPO's performance on larger-scale models and more complex RLHF challenges, including summarization and dialog tasks. Our findings indicate that, with minimal hyperparameter tuning, DPO often matches or surpasses the performance of competitive methods like RLHF with PPO, and can also outperform approaches that select the top trajectory from $N$ candidates using a learned reward function. Prior to presenting the empirical outcomes, we outline the experimental protocol; further specifics can be found in Appendix~\ref{app:exp_details}.

\textbf{Tasks.} Our study investigates three distinct tasks within the domain of open-ended text generation. In every case, learning is carried out by training a policy on a dataset of preferences $\mathcal{D}=\bigl\{x^{(i)}, y_w^{(i)}, y_l^{(i)}\bigr\}_{i=1}^N$. For the \textbf{controlled sentiment generation} task, $x$ represents the beginning segment of a movie review taken from the IMDb dataset \cite{maas-EtAl:2011:ACL-HLT2011}, and the objective is for the policy to produce $y$ exhibiting positive sentiment. To enable systematic comparison, preference pairs between generated outputs are created using a pre-trained sentiment classifier, such that $p(\text{positive}\mid x,y_w)>p(\text{positive}\mid x,y_l)$. The SFT approach involves continued training of GPT-2-large until it converges using movie reviews from the training set of the IMDB dataset (see App~\ref{app:sentiment_details} for further explanation). In the \textbf{summarization} scenario, $x$ refers to a Reddit forum post, and the task requires generating a summary $y$ capturing the post’s main arguments. Consistent with prior research, we employ the Reddit TL;DR summarization dataset \citep{volske-etal-2017-tl} and human preference data collected by \citeauthor{stiennon2022learning}. For SFT, we utilize a model refined on summaries authored by humans,\footnote{\url{https://huggingface.co/CarperAI/openai_summarize_tldr_sft}} leveraging the TRLX \citep{leandro_von_werra_2023_7790115} system to facilitate RLHF. Note that the human preference data was curated from outputs of a separate, comparably-trained SFT model as per \citeauthor{stiennon2022learning}. Lastly, in the \textbf{single-turn dialogue} context, $x$ denotes a user input that ranges from factual inquiries (e.g., astrophysics) to personal questions. The task here is for the policy to craft a response $y$ that is both engaging and informative; we use the Anthropic Helpful and Harmless dialogue dataset \citep{bai2022training}, comprising 170,000 conversations between a human participant and an AI assistant. Each conversation culminates in two candidate model-generated responses, with human annotators identifying their preference. Since a pre-trained SFT model is absent in this task, we build the SFT model by fine-tuning a pre-existing language model exclusively on responses labeled as preferred.

\textbf{Evaluation.} We assess the performance of each algorithm using two distinct evaluation methodologies. To specifically measure how well algorithms achieve the goal of constrained reward maximization, we conduct controlled sentiment generation experiments and report the Pareto frontier defined by achieved reward and KL-divergence relative to the reference policy. This is feasible because the underlying reward signal—provided by a sentiment classifier—is known in this setup. However, in practical scenarios where the exact reward function is unavailable, we rely on the \textit{win rate} metric, comparing each model to a baseline policy. For this, GPT-4 serves as a surrogate for human judgment, evaluating the comparative quality of summaries and dialogue responses in the summarization and single-turn dialogue tasks, respectively. The reference summary in the test set acts as the baseline in summarization, whereas for dialogue, we use the annotated preferred response. While prior research indicates language models can outperform standard metrics as evaluators \citep{Chen2023ExploringTU}, we supplement our approach with a human evaluation study to further validate using GPT-4 as an evaluation proxy (see Sec.~\ref{sec:human-judgments}). Results demonstrate that GPT-4 assessments strongly align with human preferences, showing human-GPT-4 agreement at least on par with, if not exceeding, human-human agreement.

\textbf{Methods.} Alongside DPO, our study includes multiple established methods for training language models according to human preferences. For the summarization task, we test zero-shot prompting using \textbf{GPT-J} \citep{gpt-j}, and for dialogue, 2-shot prompting with \textbf{Pythia-2.8B} \citep{biderman2023pythia}. We also evaluate the \textbf{SFT} model, as well as \textbf{Preferred-FT}, which is trained via supervised fine-tuning on the preferred completion $y_w$, using outputs from SFT (in controlled sentiment and summarization) or from a generic language model (in dialogue). The \textbf{Unlikelihood} approach~\citep{welleck2019neural} is another pseudo-supervised technique that increases the likelihood of $y_w$ while decreasing that of $y_l$, applying an optional scaling factor $\alpha\in[0,1]$ to the unlikelihood penalty. Furthermore, we examine \textbf{PPO} \citep{schulman2017proximal} with a reward model learned from preference data, as well as \textbf{PPO-GT}, which has oracle access to the ground truth reward function in the sentiment experiments. For PPO-GT, we employ both an out-of-the-box version \cite{leandro_von_werra_2023_7790115} and a customized version that normalizes the reward signal and applies tuned hyperparameters for enhanced performance (with similar modifications made to standard PPO when using learned rewards). Lastly, we include the \textbf{Best of $N$} method, in which $N$ outputs are sampled from the SFT model (or from Preferred-FT in dialogue), and the completion receiving the highest reward—as predicted by the preference-trained model—is chosen. Although effective, this approach is not practical for larger values of $N$, as it requires generating $N$ candidate responses per query at test time, thus significantly increasing computational overhead.

\subsection{How well can DPO optimize the RLHF objective?}

The KL-constrained reward maximization framework employed in standard RLHF methodologies is designed to find a balance between maximizing reward and constraining the learned policy's divergence from a reference policy. Hence, when evaluating different methods, it is essential to jointly consider both the magnitude of reward obtained and the corresponding KL divergence; achieving marginally higher reward is insufficient if it is accompanied by a substantially increased KL divergence. Figure~\ref{fig:frontier-tldr-main} presents the reward versus KL divergence tradeoff curves (frontier) for several methods within the sentiment task. We perform repeated experiments for every method, each time varying the hyperparameter governing policy conservativeness—specifically, using target KL values $\in\{3,6,9,12\}$ for PPO, $\beta \in \{0.05,0.1,1,5\}$ and $\alpha\in\{0.05,0.1,0.5,1\}$ for unlikelihood training, and distinct random seeds for preferred-FT. This grid results in a total of 22 experimental runs. Following each set of 100 optimization steps, up to convergence, each resulting policy is evaluated over a designated set of held-out prompts. We then compute both the mean reward under the actual reward function and the mean sequence-level KL\footnote{Here, sequence-level KL is defined as the aggregate of KL divergences calculated at each step in the sequence.} to the reference policy, $\text{KL}\left(\pi\mid \mid \pi_\text{ref}\right)$. Our analysis indicates that DPO constructs a significantly more favorable reward/KL frontier, outperforming other methods by yielding higher rewards at consistently low KL values. This outcome is especially noteworthy for several reasons. To begin, although both DPO and PPO optimize the identical objective, DPO attains markedly better efficiency; across all tested settings, DPO’s tradeoff curve strictly dominates that of PPO. Additionally, DPO even surpasses PPO on the reward/KL frontier in settings where PPO is permitted direct access to ground truth reward signals (denoted PPO-GT).

\subsection{Can DPO scale to real preference datasets?}
\label{sec:dpo-real-datasets}
We proceed to assess the ability of DPO to fine-tune on tasks involving summarization and single-turn dialogue using datasets grounded in real human preferences. In summarization, conventional automated metrics such as ROUGE are often misaligned with human judgment~\citep{stiennon2022learning}, and previous studies have demonstrated that optimizing LMs with PPO based on human preference feedback leads to more effective summarizations. For our analysis, we draw completions from the test set of the TL;DR summarization dataset, then compute the mean win rate for each method against the ground truth completions. Outputs for all compared methods are generated at temperature settings ranging from 0.0 to 1.0, with the corresponding win rates depicted in Figure~\ref{fig:frontier-tldr-main} (right). All approaches, including DPO, PPO, and Preferred-FT, conduct fine-tuning on the same GPT-J SFT model\footnote{\url{https://huggingface.co/CarperAI/openai_summarize_tldr_sft}}. Empirically, DPO attains a win rate close to 61\% at temperature 0.0, outperforming PPO, which achieves approximately 57\% at its best (also at temperature 0.0). Furthermore, DPO secures a higher peak win rate relative to the best-of-$N$ baseline. It should be noted that the $\beta$ hyperparameter for DPO was not thoroughly optimized; thus, our reported results could be conservative with respect to DPO’s full capabilities. In addition, DPO demonstrates greater resilience to variations in the sampling temperature compared to PPO, whose performance can fall back to the baseline GPT-J model at elevated temperatures. The Preferred-FT method offers minimal gains over the initial SFT model. For a direct comparison between DPO and PPO based on human assessments, see Section~\ref{sec:human-judgments}, where DPO completions sampled at temperature 0.25 are preferred 58\% of the time relative to PPO samples generated at the same temperature.

For the single-turn dialogue setting, we conduct our evaluations using the subset of the Anthropic HH test split \citep{bai2022training} containing instances with just one round of human-assistant interaction. In our GPT-4 assessments, we use the human-preferred responses in the test data as references to determine the win rate across different approaches. Since no standard SFT model exists for this particular task, our procedure begins with the pre-trained Pythia-2.8B; we first employ Preferred-FT to fine-tune a reference model on the selected completions, ensuring outputs remain within the model's data distribution, followed by DPO training. Comparisons are made with both the top response out of 128 Preferred-FT completions (noting that performance of the Best of $N$ baseline saturates at $N=128$ for this scenario; refer to Appendix Figure~\ref{fig:best-of-n}) and a 2-shot prompted variant of the base Pythia-2.8B. We observe that, for optimal temperature settings, DPO consistently matches or exceeds the performance of these baselines. We further include an RLHF model, trained via PPO on the Anthropic HH dataset \footnote{\url{https://huggingface.co/reciprocate/ppo_hh_pythia-6B}} from a reputable implementation \footnote{\url{https://github.com/CarperAI/trlx/tree/main/examples/hh}}, but are unable to identify prompting or sampling temperatures that surpass the base Pythia-2.8B’s performance. Given the parallels between the reward functions optimized by PPO and Best of 128, as established by our TL;DR experiments, we regard Best of 128 as an approximate proxy for PPO outcomes. In summary, DPO emerges as the only approach that achieves computational efficiency while reliably improving over the preferred completions in the Anthropic HH dataset, delivering performance on par with or superior to the computationally intensive Best of 128. Additionally, Figure~\ref{fig:dialogue-main} indicates that DPO attains peak performance with relatively few training iterations.

\subsection{Generalization to a new input distribution}

To further analyze how PPO and DPO handle distributional shifts, we test the policies trained on Reddit TL;DR summarization using news articles drawn from the test split of the CNN/DailyMail dataset \citep{nallapati-etal-2016-abstractive}. The policies are evaluated using the optimal sampling temperatures identified for TL;DR, specifically 0 and 0.25. Table~\ref{tab:ood} summarizes the outcomes. We measured the win rate of GPT-4 when comparing generated summaries to the ground-truth references, adapting our previously used GPT-4 (C) prompt for Reddit TL;DR by swapping the phrase ``forum post'' for ``news article''. Across this new input domain, DPO maintains a considerable advantage over PPO in terms of performance. These findings preliminarily suggest that DPO is capable of generalizing at least as effectively as PPO, even though it does not make use of the extra unlabeled Reddit TL;DR prompts utilized in PPO training.

\subsection{Validating GPT-4 judgments with human judgments}
\label{sec:human-judgments}

We undertake a human evaluation study to assess the dependability of GPT-4's preference judgments, employing outputs from the TL;DR summarization tasks and two GPT-4 prompt variants. The \textbf{GPT-4 (S)} (simple) prompt solicits a direct comparison on which summary better encapsulates the salient information from the original post. The \textbf{GPT-4 (C)} (concise) prompt also requests a judgment about conciseness, as we observe that GPT-4, under the simple prompt, tends to favor overly lengthy and repetitive summaries relative to human raters. Detailed prompt wording can be found in Appendix~\ref{app:prompts}. Our study spans three comparison settings: one using the best-performing method (DPO, temperature 0.25), one with the lowest-performing (PPO, temperature 1.0), and one with a middle-range model (SFT, temperature 0.25), all evaluated against the PPO method that uses greedy sampling (the top-performing temperature). This design aims to encompass a range of summary qualities. In both prompt conditions, the concordance between GPT-4's judgments and human ratings closely matches inter-human agreement rates, which suggests that GPT-4 offers a viable stand-in for human evaluation (with the caveat that only the DPO and PPO-1 comparisons included repeated human raters due to participant limits). Notably, the \textbf{GPT-4 (C)} prompt generally yields results more closely aligned with human preferences, so it is adopted as the primary evaluation metric in Section~\ref{sec:dpo-real-datasets}. Comprehensive information about the human study, including the interface shown to raters and the identities of the participants, can be found in Appendix~\ref{app:human-study}.

\section{Discussion}
Preference-based learning offers a robust and adaptable approach for developing proficient and aligned language models. In this work, we have presented DPO, a streamlined training scheme that leverages human preferences without reliance on reinforcement learning methodologies. Instead of adapting the preference learning challenge into a conventional RL framework to utilize standard RL techniques, DPO establishes a direct relationship between the policies of language models and reward structures. This connection makes it possible to train models to adhere to human preferences using a straightforward cross-entropy loss, fully circumventing the need for reinforcement learning while maintaining general applicability. Remarkably, DPO achieves performance that matches or surpasses current RLHF methods—such as those utilizing PPO—while requiring almost no hyperparameter tuning. As a result, DPO substantially lowers the practical obstacles involved in preference-based training of language models.

\textbf{Limitations \& Future Work.} Several pressing avenues for subsequent investigation arise from our findings. A key question is the out-of-distribution generalization ability of policies trained using DPO compared to those trained via explicit reward maximization; initial experiments indicate that DPO-trained policies generalize comparably to models optimized with PPO, though extensive examination is warranted. Furthermore, there is the open question of whether self-labeling with the DPO policy can capitalize on unlabeled prompts as effectively. Another topic concerns how over-optimization of the implicit reward is manifested in DPO, and whether the minor dip in performance seen in Figure~\ref{fig:dialogue-main}-right reflects such effects. Although our study assesses models up to a 6B parameter scale, advancing DPO to much larger, state-of-the-art models is a promising area for future research. In terms of assessment, our findings reveal that GPT-4-calculated win rates are sensitive to prompting, suggesting further inquiry into methods for reliably obtaining high-fidelity evaluations from automated systems. Lastly, DPO’s potential extends well beyond its application to language models trained on human preference signals; it could also be leveraged for optimizing generative models across a range of other data modalities.